% GENERATED FILE. Edit canonical lessons, then run npm run generate:books.
% canonical-source-hash: fnv1a64:7a741b7c7fa2839a
% canonical-lessons: TE-C342-dandora, TE-C342-pramukhudu, TE-C342-sabha, TE-C342-prasangam, TE-C342-prasangincu

\chapter{A Public Meeting}
\label{ch:te-a-public-meeting}
\hlchaptermodality{342}

\begin{chapteropening}
\textbf{By the end of this chapter:} \emph{I can say a drum proclamation, a notable, a public meeting, a speech and to give a speech.}

Complete the last lesson of chapter 342: I can say a drum proclamation, a notable, a public meeting, a speech and to give a speech.
\end{chapteropening}

\section[\texorpdfstring{\te{దండోరా} (daṇḍōrā)}{దండోరా (daṇḍōrā)}]{\te{దండోరా} (daṇḍōrā) — a drum proclamation}

\label{lesson:TE-C342-dandora}



\begin{quote}
Before the new one: say the Telugu for a reader, then the Telugu for to certify.
\end{quote}

\subsection*{You'll want to know: \te{దండోరా}}
\textbf{\te{దండోరా}} — \emph{daṇḍōrā} — \textquotedblleft{}a drum proclamation\textquotedblright{}.

\begin{quote}
\te{గ్రామంలో} \te{ఒక} \te{దండోరా} \te{వినిపించింది}.
\end{quote}

\begin{quote}
\emph{grāmaṁlō oka daṇḍōrā vinipiñcindi.}
\end{quote}

\begin{quote}
A drum proclamation was heard in the village.
\end{quote}

\te{దండోరా} \te{వేయు}, to proclaim with a drum: the old way to spread news through the whole \te{ఊరు}.

\begin{cousinweb}
\textbf{\te{దండోరా}} (\emph{daṇḍōrā}) is an old word for news cried with a drum through the streets; its origin is not settled.
\end{cousinweb}

\begin{grammarlens}[title={What you've built}]
One more word for telling what happened.
\end{grammarlens}

\subsection*{Your turn}
\begin{itemize}
\raggedright
  \item \emph{Say it:} \emph{daṇḍōrā}
  \item \emph{Say it:} \emph{daṇḍōrā}, once more
  \item \emph{Say it:} \emph{grāmaṁlō oka daṇḍōrā vinipiñcindi.} — A drum proclamation was heard in the village.
  \item \emph{Recall:} say the Telugu for a reader, then the Telugu for to certify, then say \emph{daṇḍōrā} again
\end{itemize}

\begin{tcolorbox}[breakable,colback=teal!4,colframe=teal!35!black,title={Before you move on}]
What does \te{దండోరా} mean? (\textquotedblleft{}A drum proclamation\textquotedblright{}.) Say it once more.
\end{tcolorbox}
\section[\texorpdfstring{\te{ప్రముఖుడు} (pramukhuḍu)}{ప్రముఖుడు (pramukhuḍu)}]{\te{ప్రముఖుడు} (pramukhuḍu) — a well-known person; a notable}

\label{lesson:TE-C342-pramukhudu}



\begin{quote}
Before the new one: say the Telugu for to certify, then the Telugu for a drum proclamation.
\end{quote}

\subsection*{You'll want to know: \te{ప్రముఖుడు}}
\textbf{\te{ప్రముఖుడు}} — \emph{pramukhuḍu} — \textquotedblleft{}a well-known person; a notable\textquotedblright{}.

\begin{quote}
\te{ఆ} \te{పెళ్ళికి} \te{చాలా} \te{మంది} \te{ప్రముఖులు} \te{వచ్చారు}.
\end{quote}

\begin{quote}
\emph{ā peḷḷiki cālā mandi pramukhulu vaccāru.}
\end{quote}

\begin{quote}
Many well-known people came to that wedding.
\end{quote}

Plural: \te{ప్రముఖులు}, \emph{pramukhulu}.

\te{ప్రముఖ} alone works as an adjective: \te{ప్రముఖ} \te{రచయిత}, a well-known writer.

\begin{cousinweb}
\textbf{\te{ప్రముఖుడు}} (\emph{pramukhuḍu}) is built on Sanskrit \emph{pramukha}, \textquotedblleft{}facing forward; foremost, chief\textquotedblright{} (Monier-Williams), with the Telugu ending \emph{-uḍu}.
\end{cousinweb}

\begin{grammarlens}[title={What you've built}]
One more word for telling what happened.
\end{grammarlens}

\subsection*{Your turn}
\begin{itemize}
\raggedright
  \item \emph{Say it:} \emph{pramukhuḍu}
  \item \emph{Say it:} \emph{pramukhuḍu}, once more
  \item \emph{Say it:} \emph{ā peḷḷiki cālā mandi pramukhulu vaccāru.} — Many well-known people came to that wedding.
  \item \emph{Recall:} say the Telugu for to certify, then the Telugu for a drum proclamation, then say \emph{pramukhuḍu} again
\end{itemize}

\begin{tcolorbox}[breakable,colback=teal!4,colframe=teal!35!black,title={Before you move on}]
What does \te{ప్రముఖుడు} mean? (\textquotedblleft{}A well-known person; a notable\textquotedblright{}.) Say it once more.
\end{tcolorbox}
\section[\texorpdfstring{\te{సభ} (sabha)}{సభ (sabha)}]{\te{సభ} (sabha) — a public meeting; an assembly}

\label{lesson:TE-C342-sabha}



\begin{quote}
Before the new one: say the Telugu for a drum proclamation, then the Telugu for a well-known person.
\end{quote}

\subsection*{You'll want to know: \te{సభ}}
\textbf{\te{సభ}} — \emph{sabha} — \textquotedblleft{}a public meeting; an assembly\textquotedblright{}.

\begin{quote}
\te{సాయంత్రం} \te{గుడి} \te{దగ్గర} \te{సభ} \te{జరిగింది}.
\end{quote}

\begin{quote}
\emph{sāyantraṁ guḍi daggara sabha jarigindi.}
\end{quote}

\begin{quote}
In the evening a meeting was held near the temple.
\end{quote}

Plural: \te{సభలు}, \emph{sabhalu}.

\te{సభ} \te{జరిగింది}, a meeting took place.

\begin{cousinweb}
\textbf{\te{సభ}} (\emph{sabha}) is borrowed from Sanskrit \emph{sabhā}, \textquotedblleft{}an assembly, a hall\textquotedblright{} (CDIAL 13167).
\end{cousinweb}

\begin{grammarlens}[title={What you've built}]
One more word for telling what happened.
\end{grammarlens}

\subsection*{Your turn}
\begin{itemize}
\raggedright
  \item \emph{Say it:} \emph{sabha}
  \item \emph{Say it:} \emph{sabha}, once more
  \item \emph{Say it:} \emph{sāyantraṁ guḍi daggara sabha jarigindi.} — In the evening a meeting was held near the temple.
  \item \emph{Recall:} say the Telugu for a drum proclamation, then the Telugu for a well-known person, then say \emph{sabha} again
\end{itemize}

\begin{tcolorbox}[breakable,colback=teal!4,colframe=teal!35!black,title={Before you move on}]
What does \te{సభ} mean? (\textquotedblleft{}A public meeting; an assembly\textquotedblright{}.) Say it once more.
\end{tcolorbox}
\section[\texorpdfstring{\te{ప్రసంగం} (prasaṅgaṁ)}{ప్రసంగం (prasaṅgaṁ)}]{\te{ప్రసంగం} (prasaṅgaṁ) — a speech}

\label{lesson:TE-C342-prasangam}



\begin{quote}
Before the new one: say the Telugu for a well-known person, then the Telugu for a public meeting.
\end{quote}

\subsection*{You'll want to know: \te{ప్రసంగం}}
\textbf{\te{ప్రసంగం}} — \emph{prasaṅgaṁ} — \textquotedblleft{}a speech\textquotedblright{}.

\begin{quote}
\te{సభలో} \te{అతని} \te{ప్రసంగం} \te{చాలా} \te{బాగుంది}.
\end{quote}

\begin{quote}
\emph{sabhalō atani prasaṅgaṁ cālā bāgundi.}
\end{quote}

\begin{quote}
His speech at the meeting was very good.
\end{quote}

Plural: \te{ప్రసంగాలు}, \emph{prasaṅgālu}.

\begin{cousinweb}
\textbf{\te{ప్రసంగం}} (\emph{prasaṅgaṁ}) is borrowed from Sanskrit \emph{prasaṅga} (Monier-Williams), a word whose first sense is \textquotedblleft{}attachment, connection\textquotedblright{}; in Telugu it means a public talk.
\end{cousinweb}

\begin{grammarlens}[title={What you've built}]
One more word for telling what happened.
\end{grammarlens}

\subsection*{Your turn}
\begin{itemize}
\raggedright
  \item \emph{Say it:} \emph{prasaṅgaṁ}
  \item \emph{Say it:} \emph{prasaṅgaṁ}, once more
  \item \emph{Say it:} \emph{sabhalō atani prasaṅgaṁ cālā bāgundi.} — His speech at the meeting was very good.
  \item \emph{Recall:} say the Telugu for a well-known person, then the Telugu for a public meeting, then say \emph{prasaṅgaṁ} again
\end{itemize}

\begin{tcolorbox}[breakable,colback=teal!4,colframe=teal!35!black,title={Before you move on}]
What does \te{ప్రసంగం} mean? (\textquotedblleft{}A speech\textquotedblright{}.) Say it once more.
\end{tcolorbox}
\section[\texorpdfstring{\te{ప్రసంగించు} (prasaṅgiñcu)}{ప్రసంగించు (prasaṅgiñcu)}]{\te{ప్రసంగించు} (prasaṅgiñcu) — to give a speech}

\label{lesson:TE-C342-prasangincu}



\begin{quote}
Before the new one: say the Telugu for a public meeting, then the Telugu for a speech.
\end{quote}

\subsection*{You'll want to know: \te{ప్రసంగించు}}
\textbf{\te{ప్రసంగించు}} — \emph{prasaṅgiñcu} — \textquotedblleft{}to give a speech\textquotedblright{}.

\begin{quote}
\te{మా} \te{మాస్టారు} \te{సభలో} \te{ప్రసంగించారు}.
\end{quote}

\begin{quote}
\emph{mā māsṭāru sabhalō prasaṅgiñcāru.}
\end{quote}

\begin{quote}
Our teacher gave a speech at the meeting.
\end{quote}

Past (I …): \te{ప్రసంగించాను}, \emph{prasaṅgiñcānu}.

The verb of \te{ప్రసంగం}, made with -\te{ఇంచు}.

\begin{cousinweb}
\textbf{\te{ప్రసంగించు}} (\emph{prasaṅgiñcu}) is built on Sanskrit \emph{prasaṅga} (Monier-Williams), the source of \te{ప్రసంగం}, a speech, with the Telugu verb ending \emph{-iñcu}.
\end{cousinweb}

\begin{grammarlens}[title={What you've built}]
One more word for telling what happened.
\end{grammarlens}

\subsection*{Your turn}
\begin{itemize}
\raggedright
  \item \emph{Say it:} \emph{prasaṅgiñcu}
  \item \emph{Say it:} \emph{prasaṅgiñcu}, once more
  \item \emph{Say it:} \emph{mā māsṭāru sabhalō prasaṅgiñcāru.} — Our teacher gave a speech at the meeting.
  \item \emph{Recall:} say the Telugu for a public meeting, then the Telugu for a speech, then say \emph{prasaṅgiñcu} again
\end{itemize}

\begin{tcolorbox}[breakable,colback=teal!4,colframe=teal!35!black,title={Before you move on}]
What does \te{ప్రసంగించు} mean? (\textquotedblleft{}To give a speech\textquotedblright{}.) Say it once more.
\end{tcolorbox}
